\documentclass{article}

\IfFileExists{styles/main.sty}{%
  \usepackage[vancouver]{styles/main}%
}{%
  \usepackage[vancouver]{../../styles/main}%
}

\usepackage{xurl}
\usepackage{cleveref}

\IfFileExists{text.bib}{%
  \addbibresource{text.bib}%
}{%
  \addbibresource{pws/2/text.bib}%
}

\IfFileExists{figures/}{%
  \graphicspath{{./figures/}}%
}{%
  \graphicspath{{./pws/2/figures/}}%
}

\begin{document}

\selectlanguage{ukrainian}

\crefname{figure}{рис.}{рис.}
\Crefname{figure}{Рисунок}{Рисунки}
\crefname{table}{табл.}{табл.}
\Crefname{table}{Таблиця}{Таблиці}
\crefname{section}{розд.}{розд.}
\Crefname{section}{Розділ}{Розділи}
\crefname{equation}{формула}{формули}
\Crefname{equation}{Формула}{Формули}
\renewcommand{\crefrangeconjunction}{~--~}
\renewcommand{\crefpairconjunction}{ та~}
\renewcommand{\crefmiddleconjunction}{, }
\renewcommand{\creflastconjunction}{ та~}

\titleformat{\section}[block]
{\filcenter\bfseries\fontsize{16pt}{24pt}\selectfont}
{}
{0pt}
{\MakeUppercase}
\titlespacing*{\section}{0pt}{20pt}{6pt}

\setcounter{section}{1}
\setcounter{subsection}{1}
\subsection{Порядок виконання роботи}

\subsubsection{Текст до теми БКР «Розробка системи керування рухом нижніх кінцівок двоногого робота»}

Керування рухом нижніх кінцівок двоногого робота ускладнюється нелінійною динамікою, розривними контактами з поверхнею та обмеженнями на крутні моменти приводів. Для розв'язання задачі динамічної рівноваги застосовують модельно-прогнозувальне керування (MPC) на основі критерію точки нульового моменту (ZMP)~\cite{xuZMPAwareTaskFormulation2026}, що оптимізує контактні сили на скінченному горизонті. Перенесення отриманих сил на суглоби здійснює контролер усього тіла (WBC) із контролем допустимості контактів~\cite{zhangFeasibilityawareWholebodyControl2026}, який запобігає ковзанню стоп і стабілізує позу під час руху складним рельєфом.

Для підвищення робастності до зовнішніх збурень модельно-орієнтовані алгоритми доповнюють глибоким навчанням з підкріпленням (DRL). Автоматичне формування функцій винагороди суттєво спрощує перенесення політик із симуляції на фізичного робота (sim-to-real)~\cite{luDeepReinforcementLearning2026}. Своєю чергою, злиття мультимодальних сенсорних потоків (зору та пропріоцепції)~\cite{jiLearningbasedLocomotionControl2025} і використання перцептивних внутрішніх моделей~\cite{longLearningHumanoidLocomotion2025} дозволяють адаптувати параметри кроку до невідомих перешкод та геометрії опорної поверхні в реальному часі.

Критичним аспектом практичної реалізації системи керування є надання формальних гарантій безпеки для уникнення падінь робота. Використання ризик-чутливих безакторних стратегій на основі опуклої оптимізації~\cite{zhangRisksensitiveActorfreePolicy2023} дає змогу обмежити ймовірність виходу за межі допустимих станів. Комплексна інтеграція MPC, WBC та безпечного машинного навчання формує цілісну й енергоефективну архітектуру керування нижніми кінцівками двоногого робота.

\clearpage
\nocite{*}
\printbibliography[heading=bibintoc,title={СПИСОК ВИКОРИСТАНИХ ДЖЕРЕЛ}]

\end{document}
